% Standalone problem document, generated from the adjacent README.md.
% Compile with XeLaTeX. Mathematical content is maintained in Markdown.
\documentclass[11pt,a4paper]{article}
\usepackage[fontsize=10.5pt]{scrextend}
\usepackage[margin=22mm,headheight=15pt,headsep=9mm,footskip=11mm]{geometry}
\usepackage{fontspec}
\setmainfont{texgyrepagella}[Extension=.otf,UprightFont=*-regular,BoldFont=*-bold,ItalicFont=*-italic,BoldItalicFont=*-bolditalic]
\setsansfont{texgyreheros}[Extension=.otf,UprightFont=*-regular,BoldFont=*-bold,ItalicFont=*-italic,BoldItalicFont=*-bolditalic]
\setmonofont{texgyrecursor}[Extension=.otf,UprightFont=*-regular,BoldFont=*-bold,ItalicFont=*-italic,BoldItalicFont=*-bolditalic,Scale=MatchLowercase]
\usepackage{amsmath,amssymb}
\usepackage{unicode-math}
\setmathfont{texgyrepagella-math.otf}
\usepackage{xcolor,microtype,enumitem,needspace,fancyhdr}
\usepackage{bookmark}
\definecolor{ink}{HTML}{17344B}
\definecolor{accent}{HTML}{197D80}
\definecolor{muted}{HTML}{596773}
\hypersetup{colorlinks=true,linkcolor=accent,urlcolor=accent,pdftitle={MI-22 - Lemos--Soares
singular-value
log-majorization},pdfauthor={Open Problems in Numerical Linear Algebra}}
\urlstyle{same}
\setlength{\parindent}{0pt}
\setlength{\parskip}{5pt plus 1pt minus 1pt}
\linespread{1.04}
\setlength{\emergencystretch}{3em}
\widowpenalty=10000
\clubpenalty=10000
\displaywidowpenalty=10000
\allowdisplaybreaks[1]
\setlist{nosep,leftmargin=1.5em}
\providecommand{\tightlist}{\setlength{\itemsep}{0pt}\setlength{\parskip}{0pt}}
\providecommand{\pandocbounded}[1]{#1}
\pagestyle{fancy}
\fancyhf{}
\fancyhead[L]{\sffamily\footnotesize\color{muted}OPEN PROBLEMS IN NUMERICAL LINEAR ALGEBRA}
\fancyhead[R]{\sffamily\bfseries\footnotesize\color{accent}MI-22}
\fancyfoot[L]{\parbox[b]{0.9\textwidth}{\sffamily\fontsize{8}{10}\selectfont\color{muted}Editorial ratings; a bounded search cannot certify that a problem remains unsolved.\\Literature check: 2026-09-11}}
\fancyfoot[R]{\sffamily\footnotesize\color{muted}\thepage}
\renewcommand{\headrulewidth}{0.3pt}
\renewcommand{\headrule}{\hbox to\headwidth{\color{accent}\leaders\hrule height \headrulewidth\hfill}}
\makeatletter
\renewcommand{\section}{\Needspace{5\baselineskip}\@startsection{section}{1}{0pt}{12pt plus 2pt minus 2pt}{4pt}{\sffamily\bfseries\large\color{ink}}}
\renewcommand{\subsection}{\Needspace{4\baselineskip}\@startsection{subsection}{2}{0pt}{9pt plus 1pt minus 1pt}{3pt}{\sffamily\bfseries\normalsize\color{ink}}}
\makeatother
\setcounter{secnumdepth}{0}
\begin{document}
{\sffamily\small\color{muted}Matrix inequalities and norms\par}
\vspace{3pt}
{\raggedright\hyphenpenalty=10000\exhyphenpenalty=10000\sffamily\bfseries\fontsize{21}{24}\selectfont\color{ink}Lemos--Soares
singular-value log-majorization\par}
\vspace{7pt}
{\sffamily\small\textbf{Difficulty:} challenging\quad\textbf{Importance:} interesting
to the community\par}
{\sffamily\small\bfseries\color{accent}Status: Solved\par}
\vspace{4pt}
{\color{accent}\hrule height 0.8pt}
\vspace{6pt}
\textbf{Rating rationale:} Controlling every partial product of singular
values is challenging despite the proved eigenvalue-modulus analogue;
the resulting norm comparisons are relevant across matrix analysis.

\subsection{Resolution --- 2026-09-11}\label{resolution-2026-09-11}

\textbf{Negative result by Matthew J. Colbrook}, Department of Applied
Mathematics and Theoretical Physics, University of Cambridge.
\textbf{Independent proof review: PASS.}

Rational positive definite \(3\times3\) matrices at \(t=1/8\) violate
the first singular-value inequality: the left operator norm exceeds
10900, while \(\|AB\|_2<10200\). Exact rational root residuals and a
proved operator-root error bound certify the actual principal powers.

The exact target is resolved. The original statement and source evidence
are retained below; its former difficulty rating is historical.

\textbf{Primary manuscript:}
\href{https://github.com/ajt60gaibb/OpenProblemsInNLA/blob/main/matrix-inequalities-and-norms/MI-22/solution.pdf}{complete
proof PDF},
\href{https://github.com/ajt60gaibb/OpenProblemsInNLA/blob/main/matrix-inequalities-and-norms/MI-22/solution.tex}{standalone
TeX}, Theorem 1.1 and its proof;
\href{https://github.com/ajt60gaibb/OpenProblemsInNLA/blob/main/matrix-inequalities-and-norms/MI-22/solution.md}{authorship
and scope}. The
\href{https://github.com/ajt60gaibb/OpenProblemsInNLA/blob/main/references/colbrook-matrix-2026-09-11/verification/reviews/MI-22-review.md}{independent
review} checks the full original argument and records its hash. The
draft was AI-assisted; this is independent agent verification, not
external human peer review or formal certification.
\href{https://github.com/ajt60gaibb/OpenProblemsInNLA/blob/main/references/colbrook-matrix-2026-09-11/README.md}{Submission
record}.

\subsection{Problem statement}\label{problem-statement}

Let \(A,B\in\mathbb C^{n\times n}\) be positive definite and
\(t\in[0,1]\), and define
\(A\#_tB=A^{1/2}(A^{-1/2}BA^{-1/2})^tA^{1/2}\). Write
\(s_1(X)\ge\cdots\ge s_n(X)\) for the singular values of \(X\).

Does \[
\prod_{j=1}^k s_j\bigl(A^t(A\#_tB)B^{1-t}\bigr)
\le\prod_{j=1}^k s_j(AB),\qquad 1\le k<n,
\] hold for every positive integer \(n\), every such \(A,B,t\), with
equality of the products for \(k=n\)? This is the assertion
\(s(A^t(A\#_tB)B^{1-t})\prec_{\log}s(AB)\). The positive semidefinite
version in the sources is recovered by continuous regularization
\(A+\varepsilon I,B+\varepsilon I\); endpoints in \(t\) can be stated
directly.

The question concerns the singular values of a product involving a
matrix geometric mean. Such multiplicative bounds imply norm and
determinant estimates useful in numerical treatment of positive definite
matrices. An analogue for eigenvalue moduli has already been proved, but
changes both sides of the comparison and does not answer this
singular-value question.

\subsection{References}\label{references}

\begin{enumerate}
\def\labelenumi{\arabic{enumi}.}
\tightlist
\item
  M. M. Ghabries, H. Abbas, B. Mourad and A. Assi, \emph{New
  log-majorization results concerning eigenvalues and singular values
  and a complement of a norm inequality}, Linear Multilinear Algebra 71
  (2023), 1228--1243; preprint Conjecture 1.1, p.~3, and §3.
  \href{https://arxiv.org/pdf/2105.13356}{Full preprint};
  \href{https://doi.org/10.1080/03081087.2022.2059050}{published
  article}.
\item
  J. Shi, Z. Wang and C. Wei, \emph{Log-majorizations of Lemos--Soares
  type for the metric geometric and spectral geometric means}, Ann.
  Funct. Anal. 17 (2026), article 53, p.~3, equation (1.2) and the
  following open-status statement.
  \href{https://doi.org/10.1007/s43034-026-00532-x}{Publisher};
  \href{https://www.scribd.com/document/1073369698/s43034-026-00532-x}{publicly
  indexed full-text mirror}.
\end{enumerate}

Status check (2026-09-10): the original arXiv record lists v1. The July
14, 2026 paper reproduces precisely this inequality as (1.2) on p.~3 and
explicitly says that it is still open. The publisher PDF was
inaccessible, but the indexed primary-paper text in the public mirror
exposed this passage and the later discussion distinguishing the new
partial results. Searches for ``Lemos Soares singular log-majorization
conjecture'', the exact later title, and 2025/2026 found no subsequent
announced full resolution. This is a bounded search, not certification
of openness.
\end{document}
